Let $X$ be a finite label set with $|X|=n$.  We write $S_X$ for the set of
all bijective words on $X$, identified with the symmetric group after
choosing an ordering of $X$.

The braid arrangement of type $A_{n-1}$ has reflecting hyperplanes
\[
        x_i=x_j \qquad (i\ne j).
\]
Its chambers are total orders of the coordinates, hence are indexed by
permutations; see, for example, the standard treatment of Coxeter groups
in~\cite{BjornerBrenti2005}.  We use the convention that a word
\[
        \sigma=(\sigma_1,\ldots,\sigma_n)
\]
records the labels from first to last in the corresponding total order.

\begin{definition}
Let $P$ be a strict partial order on $X$.  The set of linear extensions of
$P$ is
\[
        L(P)=\{\sigma\in S_X:
        \operatorname{pos}_\sigma(a)<\operatorname{pos}_\sigma(b)
        \text{ whenever } a<_P b\}.
\]
The associated poset cone is the chamber-union
\[
        C(P)=L(P).
\]
\end{definition}

Here the equality $C(P)=L(P)$ means that we identify chamber-unions with
their chamber-indexing sets.  Geometrically, the relation $a<_P b$
selects one side of the reflecting hyperplane $x_a=x_b$.

\begin{theorem}[Type-$A$ halfspace poset-cone dictionary]
\label{thm:halfspace-poset-cone}
In the Coxeter complex of type $A_{n-1}$, the nonempty chamber-unions
defined by intersections of reflecting halfspaces are exactly the poset
cones $C(P)$ for labeled posets $P$ on $X$.
\end{theorem}

\begin{proof}
A reflecting halfspace choice for $x_i=x_j$ is exactly one precedence
constraint, either $i<j$ or $j<i$.  If a system of such choices is
nonempty, then the corresponding precedence constraints are acyclic; their
transitive closure is a partial order $P$.  The chambers satisfying the
original halfspace choices are exactly the total orders extending this
partial order, namely $L(P)$.

Conversely, every poset $P$ gives an intersection of reflecting
halfspaces by imposing $a<b$ for each relation $a<_P b$.  The chambers in
that intersection are precisely the linear extensions of $P$.
\end{proof}

\begin{proposition}[Orbit dictionary]
\label{prop:orbit-dictionary}
The natural action of $S_X$ on type-$A$ chambers sends poset cones to
poset cones by relabeling.  Hence the $S_X$-orbits of poset cones are in
bijection with isomorphism classes of posets on $X$.
\end{proposition}

\begin{proof}
Applying a permutation $\tau\in S_X$ to a chamber word relabels every
entry.  Thus the linear extensions of $P$ are sent to the linear
extensions of the relabeled poset $\tau P$.  This proves equivariance.
Two labeled posets lie in the same relabeling orbit exactly when they are
isomorphic.
\end{proof}

\begin{proposition}[Connectivity]
\label{prop:linear-extension-connected}
Every poset cone $C(P)$ is connected in the chamber graph.
\end{proposition}

\begin{proof}
The chamber graph of type $A$ is the adjacent-transposition graph on
words.  The graph of linear extensions of a finite poset is connected
\cite[Sec.~3.5]{Stanley2012}:
given two extensions, one can move from one to the other by swapping
adjacent incomparable elements.  These swaps are precisely chamber-graph
adjacencies that remain inside $C(P)$.
\end{proof}

\begin{remark}
The theorem in this section is deliberately stated for
halfspace convexity.  In the $A_3$ tetrahedral setting, geometric,
geodesic, and halfspace convexity coincide for connected chamber-unions.
For arbitrary type $A_{n-1}$, that stronger equivalence is not used in
this paper unless supplied by a separate proof or reference.
\end{remark}
